\section{Does ecological significance change the rate of habituation or recovery?}
\label{sec:q6}

% Part III — Updating responses over time (Q6).
% See docs/OUTLINE.md — Q6.
% Draft 1. Style models: Magrath + Blumstein (agents/STYLE.md).

If missing a real threat is costly, responses to sounds that are more likely to signal danger should habituate more slowly or recover faster than responses to sounds that are not. The prediction is simple, but testing it is not, because a falling response across repeated playbacks can arise from at least four processes: sensory habituation, learning that the signal is unreliable (\secref{sec:q5}), a decline in the motivational significance of the stimulus, and growing expectation of when the next sound will occur (\secref{sec:q7}).
% VERIFY: RankinEtAl2009 (not full-text checked): criteria listed (stimulus specificity, spontaneous recovery, dishabituation) and their role in excluding fatigue and sensory adaptation.
Formal criteria for habituation, including stimulus specificity, spontaneous recovery after rest, and dishabituation by a different stimulus, distinguish habituation from fatigue and sensory adaptation \citep{RankinEtAl2009}. These criteria do not by themselves separate habituation from the other three processes, which can produce similar curves. Two claims also need to be kept apart: that responses to ecologically significant sounds habituate in a stimulus-specific way, and that significance changes the rate or time course of habituation relative to an acoustically matched control.

Field evidence for the first claim is strong. In Richardson's ground squirrels and yellow-bellied marmots, responses declined across repeated calls from one individual and recovered when the caller changed; in ground squirrels, a new exemplar from the same caller did not restore them \citep{Hare1998,BlumsteinDaniel2004} (\secref{sec:q5}). Both studies meet the stimulus-specificity criterion for habituation, at very different timescales: calls about 3~min apart in ground squirrels and about one every 10~s in marmots. Because both define the stimulus by caller identity, they cannot distinguish habituation to an individual's voice from learning that the individual's calls are not followed by danger, which is why identity-specific recovery bears on reliability as well (\secref{sec:q5}).

Evidence for the second claim rests mainly on a single laboratory study. \citet{EvansHammond1983a} measured orienting in water-deprived male rats as suppression of licking, comparing a recorded distress squeal with a synthetic mimic matched in burst timing. Responses to all stimuli habituated rapidly within a session. After rests of one and seven days, however, orienting recovered more to the natural squeal played at 100~dB than to the mimic, and spacing trials 24~h apart slowed habituation to that stimulus.
% VERIFY: EvansHammond1983b (not full-text checked): context generalization of habituation depends on stimulus significance.
A companion study reported that generalization of habituation across contexts also depended on stimulus significance \citep{EvansHammond1983b}.
These results support the prediction for recovery and long-interval habituation, but they concern orienting in laboratory rats, not alarm behavior in free-living animals, and the natural-call advantage appeared only at the higher of two sound levels (\secref{sec:q1}). We found no modern replication, and none with neural measures.

Significance does not guarantee persistence. Rat approach to playback of prosocial 50-kHz calls habituated even when playback predicted social interaction \citep{BerzEtAl2022}, so a meaningful and rewarded call can lose its behavioral effect with repetition. Both approach and accumbens dopamine responses to 50-kHz calls also declined rapidly across playbacks \citep{WilluhnEtAl2014}, although that design cannot separate habituation from a decline in motivation to approach.
% VERIFY: BlumsteinRecapet2009 (not full-text checked): proposed anti-habituation function not tested with repeated exposure.
Acoustic features may still slow habituation: the addition of nonlinear components to marmot alarm calls raised responsiveness, and the authors proposed that such features resist habituation \citep{BlumsteinRecapet2009}, but responses to repeated exposure were not tested.

These results allow a narrower answer than the question implies. Ecologically significant sounds habituate, and in field sciurids the decline is specific to the caller. Evidence that significance changes the trajectory of habituation is limited to recovery after rest in one laboratory paradigm, and significance did not prevent habituation of approach to an appetitive call. The functional prediction applies most directly to recovery: if missed threats are costly, the response that matters is whether a sound regains its effect after silence, not how strongly it is first expressed. Testing this requires comparing natural calls with matched controls on identical schedules and measuring recovery at several intervals (\secref{sec:conclusion}). Holding the caller constant within a series and changing it at test would separate habituation from reliability learning in the same design.

\tabref{tab:summary-q6} consolidates these claims for later condensation or figure drafting.

% Intermediate summary matrix for Q6. May later be condensed, moved to an appendix, or redrawn as a figure.
\begin{table}[tbp]
\centering
\footnotesize
\singlespacing
\caption{Working summary for Q6: habituation and recovery. Four processes can yield similar declining curves; only some have diagnostic tests in the current literature.}
\label{tab:summary-q6}
\begin{tabular}{@{}p{1.5in}p{3.7in}p{1.0in}@{}}
\toprule
\textbf{Claim or process} & \textbf{Evidence} & \textbf{Status} \\
\midrule
Stimulus-/caller-specific decline & Field sciurids: decline to one caller; recovery when caller changes & Strong \\
\addlinespace[2pt]
Significance slows habituation rate & Needs natural vs matched synthetic under identical schedules & Thin / gap \\
\addlinespace[2pt]
Significance speeds recovery & More recovery to loud natural distress squeal than mimic after rest & Thin (1 study) \\
\addlinespace[2pt]
Significance prevents habituation & Prosocial 50-kHz approach habituates even when playback predicts interaction & Strong against \\
\addlinespace[2pt]
Learned unreliability & Caller history without predators (see Q5); can mimic habituation curves & Confound \\
\addlinespace[2pt]
Motivational decline & Possible when outcomes are absent or devalued; rarely isolated & Underspecified \\
\addlinespace[2pt]
Temporal expectation & Schedule effects (see Q7); must hold exposure constant & Separate account \\
\bottomrule
\end{tabular}
\end{table}

